\documentclass[11pt, a4paper]{article}

% ── Packages ─────────────────────────────────────────────────────────────────
\usepackage[margin=2.5cm]{geometry}
\usepackage{amsmath, amssymb}
\usepackage{graphicx}
\usepackage{booktabs}
\usepackage{hyperref}
\usepackage{xcolor}
\usepackage{caption}
\usepackage{subcaption}
\usepackage{float}
\usepackage{enumitem}
\usepackage{parskip}

% ── Hyperlink colours ────────────────────────────────────────────────────────
\hypersetup{
  colorlinks=true,
  linkcolor=blue!60!black,
  urlcolor=blue!60!black,
  citecolor=blue!60!black,
}

% ── Placeholder macro ────────────────────────────────────────────────────────
% Usage: \placeholder{width}{height}{caption text}
\newcommand{\placeholder}[3]{%
  \begin{figure}[H]
    \centering
    \fbox{\rule{0pt}{#2}\parbox{#1}{\centering\textcolor{gray}{[Figure placeholder: #3]}}}
    \caption{#3}
  \end{figure}%
}

% ── Answer box ───────────────────────────────────────────────────────────────
\newcommand{\answerbox}[1][6cm]{%
  \vspace{0.3em}
  \noindent\fbox{\parbox{\linewidth}{\vspace{#1}}}
  \vspace{0.5em}
}

% ── Title ────────────────────────────────────────────────────────────────────
\title{
  \textbf{Flow Matching}\\[0.3em]
  \large COMP4680/8650: Advanced Topics in Machine Learning
}
\author{Hadis Amin - u8050066}
\date{\today}

% ─────────────────────────────────────────────────────────────────────────────
\begin{document}
\maketitle

% =============================================================================
\section{Implementation Approach and Experiment Setup}
% =============================================================================

\subsection{Implementation Approach}

In this section, we organise the code into reusable building blocks shared across all four parts:

\textbf{Model architectures.} Two MLP variants are defined:
\begin{itemize}[leftmargin=1.5em]
  \item \texttt{FlowMLP} (Parts 1--3): takes $(z_t, t)$, encodes $t$ with a
        128-dim sinusoidal embedding (DiT-style), concatenates $[z_t;\, e_t] \in \mathbb{R}^{D+128}$,
        then passes through 5 hidden layers (Linear $\to$ ReLU, 256 units) and
        a final linear layer producing a $D$-dimensional output. The output is
        interpreted as either $\hat{x}$ or $\hat{v}$ depending on the
        prediction type.
  \item \texttt{MeanFlowMLP} (Part 4): extends \texttt{FlowMLP} with a second
        sinusoidal embedding for the horizon $h = t - r$. The input becomes
        $[z_t;\, e_t;\, e_h] \in \mathbb{R}^{D+256}$.
\end{itemize}

\textbf{Forward process and conversions.} We use linear interpolation
$z_t = (1-t)x + t\varepsilon$ with $t$ clipped to $[0.1, 1]$, for numerical 
stability in the conversion formulas. Helper functions implement all 
conversions derived from
$z_t = (1-t)x + t\varepsilon$ and $v = \varepsilon - x$:
\[
  v = \varepsilon - x, \qquad x = z_t - t v, \qquad v = \frac{z_t - x}{t}.
\]
These let us train and evaluate any combination of prediction type and loss
space.

\textbf{Flow matching training (Parts 1--3).} Each training step samples a
batch of $x$, draws $\varepsilon \sim \mathcal{N}(0, I)$ and
$t \sim \text{Uniform}(\epsilon, 1)$ with $\epsilon = 0.01$, forms $z_t$, and computes the
loss in the chosen space (using the conversion helpers where the prediction
and loss spaces differ).

\textbf{Sampling.} Standard flow matching uses Euler ODE integration from
$t=1$ to $t=0$ with negative step size \[
  \Delta t = -\tfrac{1}{N}, \qquad
  z \leftarrow z + \hat{v}\,\Delta t.
\]
For $x$-prediction models we convert before stepping. MeanFlow uses the direct
update \mbox{$z_r = z_t - h \cdot \hat{u}$}, where $\hat{u}$ is the predicted mean
velocity over horizon $h$.

\textbf{MeanFlow training (Part 4).} At each step we sample $t$ and $r$ 
uniformly following the same clipping as Parts 1--3 ($\epsilon = 0.01$).
With probability \texttt{fm\_ratio} $= 0.7$ we set $r \leftarrow t$
(so $h = 0$), reducing the target to the standard flow-matching anchor
$v_{\mathrm{gt}}$; otherwise $h = t - r > 0$ and the consistency target
applies:
\[
  u_{\mathrm{tgt}} = v_{\mathrm{gt}} - h \cdot \frac{du}{dt},
  \qquad
  \frac{du}{dt} = \frac{\partial u}{\partial (z,t)}(z_t, t, 0)\cdot(v_{\mathrm{gt}},\,1),
\]
computed via \texttt{torch.func.jvp} (forward-mode autodiff).
The target $u_{\text{tgt}}$ is treated as a fixed constant during
backpropagation.
Stability: $\dfrac{du}{dt}$ is clamped element-wise to $[-10, 10]$ and the
total gradient norm is clipped to $1.0$.

\subsection{Experiment Setup}

\textbf{Datasets.} Three 2D toy distributions provided by
\texttt{src/dataloader.py}: \texttt{swiss\_roll}, \texttt{gaussians},
and \texttt{circles}. Each is multiplied by a random orthogonal projection
$P \in \mathbb{R}^{2 \times D}$ to produce $D$-dimensional samples for
$D \in \{2, 8, 32\}$. The intrinsic dimensionality is always 2; only the
ambient dimension changes. Back-projection for visualization uses
$\texttt{ds.to\_2d}(s) = s \cdot P^{\top}$.

\textbf{Common hyperparameters} (used in Parts 1--3 unless stated otherwise):
\begin{center}
\begin{tabular}{ll}
  \toprule
  Hyperparameter & Value \\
  \midrule
  Architecture       & 5-layer MLP, 256 hidden, ReLU \\
  Time embedding     & Sinusoidal, 128-dim (DiT-style) \\
  Optimiser          & Adam, $\mathrm{lr} = 10^{-3}$ \\
  Batch size         & 1024 \\
  Training steps     & 25{,}000 \\
  $t$ clipping       & $\epsilon = 0.01$ \\
  Sampling           & Euler ODE, 50 steps (default) \\
  Sample count       & 2000 per visualisation \\
  \bottomrule
\end{tabular}
\end{center}

\textbf{Per-part variations.}
\begin{itemize}[leftmargin=1.5em]
  \item \textbf{Part 1.2:} v-prediction + v-loss at $D=2$ on all three
        datasets, common hyperparameters.
  \item \textbf{Part 2.2:} full $4 \times 3 \times 3 = 36$ grid (4
        prediction/loss combos $\times$ 3 ambient dims $\{2,8,32\}$ $\times$
        3 datasets). Models saved to \texttt{models/}.
  \item \textbf{Part 3:} rescue strategies for v-prediction at $D=32$ on
        \texttt{swiss\_roll}. Variants tested: hidden size $256 \to 512$,
        training steps $25\text{K} \to 100\text{K}$.
  \item \textbf{Part 4.1:} reuse the Part 2 $x$-pred + $v$-loss model at
        $D=32$ and sample at step counts $\{1, 2, 5, 10, 20, 50, 100, 200\}$.
  \item \textbf{Part 4.2:} MeanFlow trained at $D=32$ on all three datasets
        with $x$-pred + $v$-loss, $\texttt{fm\_ratio}=0.7$, hidden 256,
        25K steps. Evaluated at step counts $\{1, 2, 5, 10\}$ against the
        Part 2 baseline.
\end{itemize}

\newpage
% =============================================================================
\section{Part 1 – Warm-up (10 marks)}
% =============================================================================

\subsection{Data Visualization (3 marks)}

\textbf{Task:} Load the \texttt{swiss\_roll}, \texttt{gaussians}, and \texttt{circles}
datasets at $D=2$ (original) and $D=32$ (projected back to 2D using
\texttt{to\_2d()}), and confirm that the structure is preserved.

\begin{figure}[H]
  \centering
  \includegraphics[width=\linewidth]{figures/part1_visualization.png}
  \caption{Data visualization at $D=2$ (original) and $D=32$ (projected back
  to 2D). Each column is one dataset; the bottom row confirms the intrinsic
  structure is preserved after projection.}
  \label{fig:part1_vis}
\end{figure}

% =============================================================================
\subsection{v-Prediction Flow Matching at $D=2$ (7 marks)}

\textbf{Hyperparameters}:
\begin{center}
\begin{tabular}{ll}
  \toprule
  Hyperparameter & Value \\
  \midrule
  Architecture       & 5-layer MLP, 256 hidden, ReLU \\
  Time embedding     & Sinusoidal, 128-dim (DiT-style) \\
  Optimiser          & Adam, $\mathrm{lr} = 10^{-3}$ \\
  Batch size         & 1024 \\
  Training steps     & 25{,}000 \\
  $t$ clipping       & $\epsilon = 0.01$ \\
  Sampling           & Euler ODE, 50 steps (default) \\
  Sample count       & 2000 per visualisation \\
  \bottomrule
\end{tabular}
\end{center}

\begin{figure}[H]
  \centering
  \includegraphics[width=\linewidth]{figures/part1_generated.png}
  \caption{Part 1.2 – Generated samples vs.\ ground truth for all three
  datasets at $D=2$ using v-prediction (v-loss).}
  \label{fig:part1_gen}
\end{figure}

% =============================================================================
\section{Part 2 – Flow Matching Parameterization (34 marks)}
% =============================================================================

\subsection{Derivation: Converting between Prediction Types (4 marks)}

Starting from the linear interpolation forward process:
\[
  z_t = (1-t)\,x + t\,\varepsilon, \qquad t \in [0,1], \quad
  \varepsilon \sim \mathcal{N}(0, I)
\]
where the velocity is defined as $v = \varepsilon - x$.

\textbf{From $x$-prediction to $v$:}
Solving for $\varepsilon$ gives $\varepsilon = \dfrac{z_t - (1-t)x}{t}$.
Substituting into $v = \varepsilon - x$:
\[
  v = \frac{z_t - (1-t)x}{t} - x
    = \frac{z_t - (1-t)x - tx}{t}
    = \boxed{\frac{z_t - x}{t}}
\]

\textbf{From $v$-prediction to $x$:}
Since $\varepsilon = v + x$, substituting into the forward process:
\[
  z_t = (1-t)x + t\,\varepsilon = (1-t)x + t(v + x) = x + tv
\]
Rearranging:
\[
  \boxed{x = z_t - t\,v}
\]

These conversions allow us to train with any combination of prediction type
($x$ or $v$) and loss space ($x$-loss or $v$-loss).

% ──────────────────────────────────────────────────────────────────────────────
\subsection{Reproducing the Paper's Finding (18 marks)}

All four prediction/loss combinations were trained on three datasets
($\texttt{swiss\_roll}$, $\texttt{gaussians}$, $\texttt{circles}$) and three
dimensions ($D \in \{2, 8, 32\}$), giving 36 experiments in total.

\begin{figure}[H]
  \centering
  \includegraphics[width=\linewidth]{figures/part2_swiss_roll.png}
  \caption{Part 2 – \texttt{swiss\_roll}: generated samples for all 4
  prediction/loss combinations across $D \in \{2,8,32\}$.}
  \label{fig:part2_swiss}
\end{figure}

\begin{figure}[H]
  \centering
  \includegraphics[width=\linewidth]{figures/part2_gaussians.png}
  \caption{Part 2 – \texttt{gaussians}.}
  \label{fig:part2_gauss}
\end{figure}

\begin{figure}[H]
  \centering
  \includegraphics[width=\linewidth]{figures/part2_circles.png}
  \caption{Part 2 – \texttt{circles}.}
  \label{fig:part2_circles}
\end{figure}

% ──────────────────────────────────────────────────────────────────────────────
\subsection{Questions}

\textbf{Q1 (3 marks).}
Which prediction type scales successfully to high ambient dimensions?
At what dimension do the other prediction types begin to fail visibly?

In high dimensions, only $x$-prediction produces good samples across all datasets.
$v$-prediction begins to fail visibly at $D=32$.

\textbf{Q2 (3 marks).}
Does the choice of loss space ($x$-loss, $v$-loss) affect which prediction
types succeed or fail? What does this tell you about what determines generation
quality?

No. The choice of loss space has only a minor effect on generation quality. When prediction 
and loss spaces differ, the effective loss is implicitly reweighted by $t^2$ or $\dfrac{1}{t^2}$. 
However, $v$-prediction fails at $D=32$ regardless of loss space, confirming
that \textbf{prediction type} is the dominant factor determining generation quality.

\textbf{Q3 (6 marks).}
Explain why the successful prediction type works at high $D$ while the others
do not. Consider the nature of each prediction target and the rank of each
target relative to the ambient dimension.

As shown in \texttt{dataloader.py}, the 32D datasets are generated by linearly projecting 2D data using a matrix $P \in \mathbb{R}^{32 \times 2}$. This means the intrinsic rank of the data is only 2, even though it lives in 32-dimensional space.

Because of this, $x$-prediction succeeds at high $D$ — its target $x$ lies on a 2D manifold, so the model only needs to learn a low-rank structure regardless of how large $D$ is. The complexity of the target does not grow with $D$.

In contrast, $v$-prediction must learn $v = \varepsilon - x$, where $\varepsilon \sim \mathcal{N}(0, I_D)$ is Gaussian noise that spreads across all $D$ dimensions. The noise term dominates and forces the model to learn a full 32D target, even though the underlying data is only 2D. This is why $v$-prediction fails as $D$ grows.

% =============================================================================
\section{Part 3 – Can We Rescue v-Prediction? (25 marks)}
% =============================================================================

From Part 2, v-prediction fails at high dimensions because the data lives on a 2D manifold
but the noise $\varepsilon \sim \mathcal{N}(0, I_D)$ is generated in all $D=32$ dimensions.
The noise therefore dominates the velocity target $v = \varepsilon - x$, making it a full-rank
$D$-dimensional function that a small model cannot represent.
Consistent with RAE Section 4.1, which states that model width must exceed the token dimension
for convergence, the proposed solution is to increase model capacity via:
\begin{enumerate}
  \item Increasing the number of hidden units (256 $\to$ 512).
  \item Increasing the number of training iterations (25K $\to$ 100K).
\end{enumerate}

\begin{figure}[H]
  \centering
  \includegraphics[width=\linewidth]{figures/part3_rescue.png}
  \caption{Part 3 – Rescue attempts for v-prediction on \texttt{swiss\_roll}
  at $D=32$.}
  \label{fig:part3}
\end{figure}

\subsection{Questions}

\textbf{Q1 (2 marks).}
Based on your experiments, is v-prediction's failure fundamental or can it be
overcome?  Do your findings support or contradict your observations from Part 2?

V-prediction's failure is \textbf{not fundamental}. As observed in Part 2, $v = \varepsilon - x$ 
has full rank $D$, forcing the model to learn a full 32-dimensional function --- unlike $x$-prediction 
whose target stays rank-2 regardless of $D$. 
% Consistent with RAE Section 4.1, increasing model 
% capacity (hidden size 512 or 100K training steps) successfully rescues v-prediction. With sufficient 
% capacity, the model can learn the full complexity of the noise-dominated velocity field in high dimensions 
% ($D=32$), confirming the failure is a capacity challenge, not an inherent limitation of the parameterization.

\textbf{Q2 (3 marks).}
What approach(es) did you try?  Compare the compute cost between your approach
and the default x-prediction setup to achieve similar sample quality at $D=32$.

Two approaches successfully rescued v-prediction on \texttt{swiss\_roll} at $D=32$: 
increasing hidden size from 256 to 512, and extending training from 25K to 100K steps. 

Both successful approaches come at a significant compute cost compared to default x-prediction. The wider model increases parameter count from $\sim$313K to $\sim$1.15M ($\approx 3.7\times$ more), while the longer run requires $4\times$ more training steps. By contrast, x-prediction achieves good quality with the default 256-hidden, 25K-step setup --- making it substantially more compute-efficient at high dimensions.

\begin{figure}[H]
  \centering
  \includegraphics[width=\linewidth]{figures/part3_q3_comparison.png}
  \caption{Part 3 Q3 – v-prediction (top, red) vs.\ x-prediction (bottom, blue) response to
  capacity changes on \texttt{swiss\_roll} at $D=32$.}
  \label{fig:part3_q3}
\end{figure}

\textbf{Q3 (3 marks).}
Compare how x-prediction and v-prediction respond to your changes.  Do they
behave the same way?  Explain why or why not, considering what each
parameterisation must learn as a function of the ambient dimension $D$.

No. x-prediction is insensitive to capacity increases: the baseline already produces clean samples,
and wider or longer training changes nothing meaningful. This is because its target $x$ lies on a
rank-2 manifold regardless of $D$, so the default model has sufficient capacity.
v-prediction, by contrast, improves significantly with more capacity — because its target
$v = \varepsilon - x$ has full rank $D$, requiring a larger model to represent the more complex function.

% More precisely, x-prediction must learn $\hat{x}(z_t, t)$ whose output always lies on the 2D data
% manifold: complexity is \emph{constant}, independent of $D$.
% v-prediction must learn $\hat{v}(z_t, t) = \varepsilon - x$; since $\varepsilon \sim \mathcal{N}(0,
% I_D)$ fills all $D$ dimensions, the output complexity \emph{grows with $D$}.
% At $D=32$, the model must represent a full 32-dimensional function even though the data is only 2D
% — which is why extra capacity is necessary for v-prediction but redundant for x-prediction.

\textbf{Q4 (7 marks).}
In practice, v-prediction is used successfully in real image generation systems
such as Stable Diffusion 3 and FLUX.  Why might the situation be different for
those models compared to these toy datasets?

Based on Q2, $v$-prediction struggles because the target $v = \varepsilon - x$ is dominated by full-rank Gaussian noise, forcing the model to learn across all $D$ dimensions even though the actual data structure is low-rank. In real models like SD3/FLUX, this is not an issue because diffusion happens in a structured latent space produced by a pretrained encoder (e.g., DINOv2). The encoder compresses the input into a semantically rich representation with much lower effective rank, so the noise added during diffusion does not span the full ambient dimension. This makes $v$-prediction tractable even at high dimensions.

% =============================================================================
\section{Part 4 – One-Step Generation (31 marks)}
% =============================================================================

\subsection{Sampling Efficiency (3 marks)}

Using the best model from Part 2 ($x$-prediction, $v$-loss, $D=32$), we
evaluate how sample quality degrades as the number of Euler integration steps
decreases ($1, 2, 5, 10, 20, 50, 100, 200$ steps) across all three datasets.

\begin{figure}[H]
  \centering
  \includegraphics[width=\linewidth]{figures/part4_efficiency_swiss_roll.png}
  \caption{Part 4.1 – Sampling efficiency on \texttt{swiss\_roll} ($D=32$,
  $x$-prediction).}
  \label{fig:part4_eff_swiss}
\end{figure}

\begin{figure}[H]
  \centering
  \includegraphics[width=\linewidth]{figures/part4_efficiency_gaussians.png}
  \caption{Part 4.1 – Sampling efficiency on \texttt{gaussians} ($D=32$,
  $x$-prediction).}
  \label{fig:part4_eff_gauss}
\end{figure}

\begin{figure}[H]
  \centering
  \includegraphics[width=\linewidth]{figures/part4_efficiency_circles.png}
  \caption{Part 4.1 – Sampling efficiency on \texttt{circles} ($D=32$,
  $x$-prediction).}
  \label{fig:part4_eff_circles}
\end{figure}

% ──────────────────────────────────────────────────────────────────────────────
\subsection{MeanFlow (9 marks)}

\textbf{Architecture.}
The MeanFlow model $u(z, t, h)$ extends the base MLP with an additional
sinusoidal embedding for the horizon $h = t - r$.  The input becomes
$[z_t;\, e_t;\, e_h] \in \mathbb{R}^{D+256}$, where $e_t, e_h \in
\mathbb{R}^{128}$ use separate embedding parameters.

\textbf{Training objective.}
We use $x$-prediction with $v$-loss (the best parameterization from Part 2):
the model $u_\theta(z, t, h)$ outputs an estimate of the clean data $\hat{x}$,
and the predicted velocity is $\hat{v} = (z_t - \hat{x})/t$. Training samples
$(t, r)$ are drawn jointly with $h = t - r$. With probability
$\texttt{fm\_ratio} = 0.7$ we set $r = t$ (so $h = 0$), giving standard
flow matching:
\[
  \mathcal{L}_{\text{FM}} = \mathbb{E}\bigl[\|\hat{v}(z_t, t, 0) - v_{gt}\|^2\bigr],
  \qquad v_{gt} = \varepsilon - x.
\]
With probability $0.3$ we use $h > 0$ and apply the MeanFlow consistency loss:
\[
  \mathcal{L}_{\text{MF}} = \mathbb{E}\bigl[\,\|\hat{v}(z_t, t, h) -
    \text{sg}\bigl(v_{gt} - h \cdot \tfrac{du}{dt}\bigr)\|^2\,\bigr]
\]
where the total derivative along the flow,
\[
  \frac{du}{dt} = \frac{\partial u}{\partial z}\cdot v_{gt}
                + \frac{\partial u}{\partial t},
\]
is computed in a single forward-mode sweep via \texttt{torch.func.jvp}, and
$\text{sg}(\cdot)$ denotes stop-gradient. Note that for $h=0$ samples the
$h \cdot du/dt$ term vanishes, so the formula naturally reduces to the
flow-matching anchor target $v_{gt}$ --- meaning a single masked formulation
covers both regimes.

\textbf{Sampling.}
\[
  z_{t-h} = z_t - h \cdot u(z_t,\, t,\, h)
\]
For single-step generation: $z_0 = z_1 - 1 \cdot u(z_1, 1, 1)$.

\begin{figure}[H]
  \centering
  \includegraphics[width=\linewidth]{figures/part4_meanflow.png}
  \caption{Part 4.2 – MeanFlow samples on all three datasets at $D=32$:
  ground truth alongside $1$, $2$, $5$-step generation. Trained with
  $x$-pred + $v$-loss, $\texttt{fm\_ratio}=0.7$, $\texttt{hidden}=256$,
  $25\,000$ steps.}
  \label{fig:part4_mf_grid}
\end{figure}

\begin{figure}[H]
  \centering
  \includegraphics[width=\linewidth]{figures/part4_comparison_swiss_roll.png}
  \caption{Part 4.2 – Side-by-side: standard flow matching ($x$-pred from
  Part 2) vs.\ MeanFlow on \texttt{swiss\_roll} ($D=32$), at $1, 2, 5$ steps.}
  \label{fig:part4_cmp_swiss}
\end{figure}

\begin{figure}[H]
  \centering
  \includegraphics[width=\linewidth]{figures/part4_comparison_gaussians.png}
  \caption{Part 4.2 – Side-by-side: standard flow matching vs.\ MeanFlow on
  \texttt{gaussians} ($D=32$), at $1, 2, 5$ steps.}
  \label{fig:part4_cmp_gauss}
\end{figure}

\begin{figure}[H]
  \centering
  \includegraphics[width=\linewidth]{figures/part4_comparison_circles.png}
  \caption{Part 4.2 – Side-by-side: standard flow matching vs.\ MeanFlow on
  \texttt{circles} ($D=32$), at $1, 2, 5$ steps.}
  \label{fig:part4_cmp_circles}
\end{figure}

% ──────────────────────────────────────────────────────────────────────────────
\subsection{Questions}

\textbf{Q1 (2 marks).}
Why did you choose this prediction type for MeanFlow?
Connect your answer to your findings from Part 2.

The $x$-prediction is chosen because its target lies on a rank-2 manifold
regardless of $D$, as established in Part 2. MeanFlow requires computing
$u_{\text{tgt}}$ via a self-referential JVP; with $x$-prediction the
model output is low-rank and well-conditioned, producing stable and
consistent targets throughout training.

\textbf{Q2 (4 marks).}
In your own words, describe the core idea behind MeanFlow.  What does the
model learn that is different from standard flow matching?  Why does this
enable one-step generation?

Standard flow matching trains a model to predict the instantaneous
velocity $v(z_t, t)$ at a single timestep $t$. Meanwhile, MeanFlow trains
the model to predict the \textit{mean velocity} $u(z_t, t, h)$ — the
average velocity over a finite horizon $h = t - r$. Because the model
directly learns how to move from $z_t$ to $z_r$ in one step via
$z_r = z_t - h \cdot u$, it does not need to integrate many small
Euler steps. Setting $t=1$, $r=0$, and $h=1$ at inference gives
one-step generation: $z_0 = z_1 - u(z_1, 1, 1)$, jumping directly
from pure noise to data in a single forward pass.

\textbf{Q3 (3 marks).}
MeanFlow splits training between $h=0$ (standard flow matching) and $h>0$
(mean velocity).  Why is the $h=0$ portion needed?

The $h=0$ portion is needed because the consistency target
$u_{\text{tgt}} = v_{\text{gt}} - h \cdot \dfrac{du}{dt}$ depends on
the model itself through the JVP. This means the target changes as the
model updates, so there is no fixed ground truth for the model to
learn from. When $h=0$, the target simply becomes $v_{\text{gt}}$,
which is a true ground truth independent of the model. This gives
the training a stable reference point, so the $h>0$ consistency
loss has something reliable to build on.


\textbf{Q4 (3 marks).}
Compare the training cost of MeanFlow to standard flow matching.  Why is
MeanFlow harder to train?  What is the computational overhead of the JVP
operation per training step?

Standard flow matching requires one forward pass to compute the
prediction, one loss computation against $v_{\text{gt}}$, and one
backward pass. MeanFlow requires an additional JVP computation per
step, which itself consists of two forward passes: one for the primal
output $u$ and one for the tangent output $\frac{du}{dt}$. The total
derivative $\frac{du}{dt} = \frac{\partial u}{\partial z} \cdot
v_{\text{gt}} + \frac{\partial u}{\partial t}$ is then used to
construct the consistency target $u_{\text{tgt}} = v_{\text{gt}} -
h \cdot \frac{du}{dt}$. The overhead of the JVP is approximately two
additional forward passes per step, making MeanFlow roughly $2\times$
more expensive per training step compared to standard flow matching.

\textbf{Q5 (7 marks).}
Compare the MeanFlow-generated samples against the ground truth across all
three datasets.  Do you observe any differences or artefacts, particularly on
the \texttt{gaussians} dataset?  Describe what you see and explain why it
occurs (or why it does not).

\texttt{swiss\_roll} is the cleanest across all step counts. Because
the distribution is continuous, the average velocity consistently
points along the spiral — there is no empty space to get lost in.

\texttt{circles} falls in between — the ring recovers well at 2--5
steps with mild scatter at 1 step, as the ring is mostly continuous
but has some empty interior.

\texttt{gaussians} shows the most artifacts, with scatter appearing
between modes at low step counts. Because the distribution is
multimodal with empty space between modes, the model learns an
average velocity that can point in the wrong direction for samples
that start between modes, causing them to land in the gaps rather
than on a mode.

% =============================================================================
\end{document}
